\section{Related Work}
\label{sec-rel}
In this section, we review the related works in two aspects, namely large-scale recommendation methods and multi-scenario multi-task recommendation methods.

\paratitle{Large-scale Recommendations.}
Inspired by the success of scaling laws in large language models (LLMs)~\cite{LLM1,LLM2}, the scaling laws of recommendation models have recently been widely studied and have achieved significant success in real-world industrial recommendation systems~\cite{wukong,RankMixer,KGBRank,HSTU,MTGR,OneTrans,hyformer}.
Typically, these methods increase the model parameter size and computational complexity by stacking various efficient feature interaction structures, thereby improving the recommendation performance.
One type of method focuses on scaling the parameters and complexity of the non-sequential feature interaction module for better performance.
For example, Wukong~\cite{wukong} scales the feature interaction module by introducing a stacked factorization machine-based architecture.
Rankmixer~\cite{RankMixer} introduces multi-head token mixing and pertoken feed-forward networks to scale recommender models, while maintaining inference latency and achieving superior performance.
FAT~\cite{KGBRank} scales recommendation models by introducing field-aware transformer with field-based interaction priors in attention.
Another typical method adopts the transformer-based structure for modeling non-sequential and sequential feature in a unified way and scaling constantly.
For instance, HSTU~\cite{HSTU} sequentializes the heterogeneous feature space and then performs sequence transduction conditioned on contextual and candidate signals.
MTGR~\cite{MTGR} further scales recommendation models based on HSTU  for acceleration and leakage avoidance.
OneTrans~\cite{OneTrans} introduces a unified transformer-based backbone that simultaneously models user-behavior sequences and feature interactions via a unified tokenizer, enabling efficient parameter scaling and superior performance.
Compared to these methods, our work further explores how large-scale recommendation models can better perform in multi-scenario and multi-task settings.

\paratitle{Multi-scenario and Multi-task Recommendations.}
In real-world industrial recommendation systems, a recommendation model usually needs to serve multiple scenarios and tasks.
Multi-scenario learning (MSL) and multi-task learning (MTL) in recommender system have been widely studied in previous works~\cite{sharedbottom,MMoE,Cross-Stitch,snr,PLE,HMoE,STAR,PEPNet,hinet,mmfi,yrr,MTL1}.
For MTL, most existing methods adopt a shared-specific architecture in which task-shared and task-specific components are combined together (\eg Sharedbottom~\cite{sharedbottom}, Cross-stitch~\cite{Cross-Stitch}, SNR~\cite{snr}).
In addition, Mixture-of-Experts (MoE) based structures have also been widely proposed, aiming to enhance the effectiveness of multi-task learning by designing various expert mechanisms (\eg MMoE~\cite{MMoE}, PLE~\cite{PLE}).
For MSL, it aims to use a unified model to learn data from multiple scenarios, thereby improving recommendation performance.
Most existing work also adopts a shared-specific architecture in which scenario-shared and scenario-specific components are combined together.
The framework used in MTL can also be applied in MSL by treating the predictions for different scenarios as distinct prediction tasks (\eg SharedBottom~\cite{sharedbottom}, MMoE~\cite{MMoE}).
In addition, some specialized multi-scenario methods have also been proposed~\cite{HMoE,STAR,PEPNet,hinet,mmfi,hc2,M2M,AdaSparse,SASS,TreeMS,SARNet,SAML}.
For example, HMoE~\cite{HMoE} proposes the hybrid of implicit and explicit MoE structure which learns the scenario and task relationships.
STAR~\cite{STAR} introduces a start topology structure where a centered network is used for all scenarios and scenario-specific network is generated for each scenario.
PEPNet~\cite{PEPNet} takes the scenario related prior information to adjust the feature embeddings and hidden units through gate mechanisms.
In contrast to these approaches which are based on the shared-specific architecture, we propose a novel multi-distribution learning framework from the perspective of tokenization, aiming to enhance the learning capability and unified modeling ability of large-scale recommendation models in multi-scenario and multi-task settings.
